\section{Higher-dimensional extension: available energy}
\label{sec:3d}

The enabling property of unbinned OmniFold is that adding an observable is
adding an input feature: the step-1/step-2 classifiers simply gain a column,
with no migration-matrix or binning penalty. There is no practical D'Agostini
IBU analogue, which makes this the primary physics motivation for the technique.
We exercise it by extending the measurement to
$d^{3}\sigma/(d\pT\,d\pPar\,d\Eavail)$ with available energy as a third axis.

The significance is clear from the measurement record: every published MINERvA
differential cross section has unfolded at most two kinematic variables
simultaneously (Fig.~\ref{fig:landscape}, Table~\ref{tab:landscape}), the
double-differential measurements all relying on binned IBU. A simultaneous
\emph{three}-observable unfold of MINERvA data is, to our knowledge,
without precedent, and is a direct consequence of the unbinned formulation.

\begin{figure}[t]
  \centering
  \includegraphics[width=0.97\textwidth]{minerva_unfolding_landscape.png}
  \caption{Dimensionality of the full set of MINERvA differential cross-section
  measurements over time (number of kinematic variables unfolded
  \emph{simultaneously}; a double-differential $d^{2}\sigma/dx\,dy$ counts as
  two), coloured by channel. All $N=30$ prior measurements use binned D'Agostini
  IBU and reach at most two variables; this work's unbinned OmniFold adds a
  third. Source: the MINERvA data-release page; full list in
  Table~\ref{tab:landscape}. Dimensionality is classified from each analysis's
  headline differential result; flux, reconstruction, generator-tuning and
  total-only entries are excluded.}
  \label{fig:landscape}
\end{figure}

{\small
\begin{longtable}{@{}clcl@{}}
  \caption{Full MINERvA differential cross-section survey: the number of
  kinematic variables unfolded simultaneously (Dim.), classified from each
  analysis's headline differential result. Channels: CC-inclusive (incl.),
  CCQE/QE-like incl.\ transverse kinematic imbalance (QE), pion production CC
  and coherent ($\pi$), deep-inelastic / kaon (DIS), electron-neutrino
  ($\nu_e$), nuclear-target ratio (ratio).}
  \label{tab:landscape}\\
  \toprule
  Year & Channel; observables & Dim. & Ref. \\
  \midrule
  \endfirsthead
  \toprule
  Year & Channel; observables & Dim. & Ref. \\
  \midrule
  \endhead
  \midrule \multicolumn{4}{r}{\emph{continued\ldots}}\\ \endfoot
  \bottomrule
  \endlastfoot
  2013 & QE; $\nu_\mu$ CCQE $\pT,\pPar$            & 2 & \cite{MINERvA:2013qe} \\
  2013 & QE; $\bar\nu_\mu$ CCQE $\pT,\pPar$        & 2 & \cite{MINERvA:2013qebar} \\
  2014 & ratio; CC on C/Fe/Pb to CH               & 1 & \cite{MINERvA:2014ratio} \\
  2014 & $\pi$; coherent $\pi^\pm$ (total \& diff.) & 1 & \cite{MINERvA:2014coh} \\
  2015 & $\pi$; CC $1\pi^+$ ($T_\pi$, $\theta_\pi$) & 1 & \cite{MINERvA:2015ydy} \\
  2015 & QE; $\mu$+p final states                 & 1 & \cite{MINERvA:2015mup} \\
  2015 & $\pi$; $\bar\nu_\mu$ single $\pi^0$       & 1 & \cite{MINERvA:2015pi0} \\
  2016 & $\nu_e$; CCQE-like on CH                  & 2 & \cite{MINERvA:2016nueqe} \\
  2016 & $\pi$; $\nu_\mu/\bar\nu_\mu$ pion prod.   & 2 & \cite{MINERvA:2016pion} \\
  2016 & DIS; partonic nuclear effects            & 2 & \cite{MINERvA:2016dis} \\
  2016 & DIS; CC $K^+$ production                  & 1 & \cite{MINERvA:2016kaon} \\
  2017 & $\pi$; CC single $\pi^0$                  & 2 & \cite{MINERvA:2017pi0} \\
  2017 & QE; nuclear-dependence (multi-target)     & 1 & \cite{MINERvA:2017Adep} \\
  2018 & QE; $\bar\nu_\mu$ QE-like $\pT,\pPar$     & 2 & \cite{MINERvA:2018hba} \\
  2018 & QE; $\mu$-p mesonless TKI                 & 2 & \cite{MINERvA:2018tcb} \\
  2018 & $\pi$; coherent $\pi^\pm$ on C            & 2 & \cite{MINERvA:2018coh} \\
  2019 & QE; quasielastic-like $\nu_\mu$           & 2 & \cite{MINERvA:2019qe} \\
  2019 & $\pi$; $\bar\nu_\mu$ single $\pi^-$       & 2 & \cite{MINERvA:2019pim} \\
  2020 & incl.; CC-incl $\langle E_\nu\rangle\!\sim\!3.5$ GeV $\pT,\pPar$ & 2 & \cite{MINERvA:2020jwd} \\
  2020 & QE; nuclear binding \& TKI               & 2 & \cite{MINERvA:2020cgh} \\
  2020 & QE; high-statistics QE-like              & 2 & \cite{MINERvA:2020qehs} \\
  2020 & $\pi$; CC $\pi^0$ nuclear effects        & 1 & \cite{MINERvA:2020pi0} \\
  2021 & incl.; CC-incl $\langle E_\nu\rangle\!\sim\!6$ GeV $\pT,\pPar$ \textbf{(target)} & 2 & \cite{MINERvA:2021owq} \\
  2022 & QE; QE-like muon \& proton kinematics    & 2 & \cite{MINERvA:2022qe} \\
  2022 & incl.; CC-incl low momentum transfer     & 1 & \cite{MINERvA:2022incl} \\
  2023 & $\pi$; CC $1\pi^+$ (multi-target)        & 2 & \cite{MINERvA:2023ccpi} \\
  2023 & $\pi$; coherent $\pi^+$ (multi-target)   & 1 & \cite{MINERvA:2023coh} \\
  2023 & QE; $\bar\nu_\mu$ multi-neutron low-$\Eavail$ & 1 & \cite{MINERvA:2023lowEavail} \\
  2024 & $\nu_e$; low-$Q^2$ ($\Eavail$, $\pT$)    & 2 & \cite{MINERvA:2023avz} \\
  2025 & QE; $A$-dependence TKI (multi-target)    & 2 & \cite{MINERvA:2025tki} \\
  \midrule
  2026 & \textbf{this work}; unbinned $\pT,\pPar$            & 2 & --- \\
  2026 & \textbf{this work}; unbinned $\pT,\pPar,\Eavail$    & \textbf{3} & --- \\
\end{longtable}
}

\subsection{Definitions and pipeline}
\label{sec:3d-method}
Available energy follows the MINERvA low-recoil convention
(Ref.~\cite{MINERvA:2023avz}, Eq.~4): truth
$\Eavail=\sum_{p}\mathrm{KE}+\sum_{\pi^{\pm}}\mathrm{KE}+\sum_{\pi^{0},\gamma}E$
(excluding neutrons and the muon), via \texttt{CVUniverse::GetEAvailableTrue()};
reco $\Eavail$ is the calorimetric tracker$+$ECAL estimate scaled by $1.17$,
\texttt{CVUniverse::NewEavail()}. Both accessors were added to the event loop,
which was re-run over all twelve ME~FHC playlists to append the $\Eavail$
branches alongside the existing $\pT$/$\pPar$ branches. The unfolding driver is
the 2D pipeline with a three-column feature stack
$(\pT,\pPar,\Eavail)$; the cross-section extraction generalises to a 3D Jacobian
$\Delta\pT\,\Delta\pPar\,\Delta\Eavail$ with the flux $\Phi$ still applied per
$\pT$ (broadcast over $\pPar$ and $\Eavail$). The $\Eavail$ binning
$[0,0.1,0.2,0.4,0.8,1.5,3.0,100]\,\si{GeV}$ ends in a wide catch bin so that the
$\Eavail$-marginal captures the full CC-inclusive recoil tail; the global
OmniFold completeness is $c=1.0000$, as in 2D.

\subsection{The available-energy spectrum}
\label{sec:3d-spectrum}
Figure~\ref{fig:eavail} shows the unfolded $d\sigma/d\Eavail$ (the 1D projection
of the 3D result). It falls monotonically from
\SI{2.19e-38}{cm^2/GeV/nucleon} at the low-recoil peak to
\SI{6.8e-41}{} in the catch bin --- the expected CC-inclusive shape, with most
of the rate at $\Eavail\lesssim\SI{0.8}{GeV}$.

\begin{figure}[t]
  \centering
  \includegraphics[width=0.66\textwidth]{eavail_spectrum.png}
  \caption{Unfolded available-energy spectrum $d\sigma/d\Eavail$, the 1D
  projection of the 3D measurement. The wide $[3,100]\,\si{GeV}$ catch bin
  (value annotated) is excluded from the axis.}
  \label{fig:eavail}
\end{figure}

\subsection{Validation against the 2D measurement}
\label{sec:3d-anchor}
With no published 3D reference, the validation anchor is that marginalising the
3D cross section over $\Eavail$ must reproduce the 2D result of
\S\ref{sec:results}. The marginal recovers the 2D \emph{normalisation} closely:
total $\sigma$ agrees to $+\SI{0.95}{\percent}$ and the per-bin
marginal/2D ratio has median $1.0016$. The \emph{shape} agreement is looser ---
the marginal versus the published covariance gives $\chi^{2}/\mathrm{ndf}=4.98$,
above the frozen 2D values ($3.66$ default, $2.65$ LightGBM-CV), driven by a
$\sim\SI{4.4}{\percent}$ per-bin scatter that the third axis injects into the
$(\pT,\pPar)$ marginal (Fig.~\ref{fig:eavail-marg}).

\begin{figure}[t]
  \centering
  \includegraphics[width=0.7\textwidth]{eavail_marginal_vs_paper_pull_full.png}
  \caption{Per-bin pull of the $\Eavail$-marginal against the published
  covariance, on the $14\times16$ grid and its distribution.}
  \label{fig:eavail-marg}
\end{figure}

\subsection{Closure: the scatter is physical, not a method bias}
\label{sec:3d-closure}
A reweight closure isolates whether that $\SI{4.4}{\percent}$ scatter is method
bias or real information. We inject a known $+\SI{30}{\percent}$ Gaussian bump
in \emph{truth} $\Eavail$ (centre \SI{0.3}{GeV}, width \SI{0.15}{GeV}) into the
pseudo-data and require the 3D unfold to recover it. It does: the recovered
step-2 mean weight is $1.048$ against the injected $1.0485$, and the residual of
the unfold against the reweighted-truth reference is, on the $\Eavail$-marginal,
median $0.9999$ with standard deviation \textbf{\num{0.0006}}
(\SI{0.06}{\percent}); the per-3D-bin residual standard deviation is
\SI{3.2}{\percent} (MC statistics in sparse high-recoil corners). The closure
marginal scatter is thus $\sim70\times$ tighter than the
$\SI{4.4}{\percent}$ data-versus-2D scatter. The method is therefore unbiased on
the new axis, and we read the elevated marginal $\chi^{2}$ as
\emph{genuine data\,$\leftrightarrow$\,MC structure that the available-energy
dimension exposes} --- information that is invisible to the 2D measurement ---
rather than a pipeline artefact (the total normalisation matches to
$<\SI{1}{\percent}$).

This establishes the 3D framework and its first closure. The full 3D
\emph{systematic}-uncertainty campaign (universe, bootstrap, and ML covariances,
as in \S\ref{sec:systematics} for 2D) is the natural follow-up; the present 3D
result is central-value only.

\paragraph{Model comparison in 3D.}
The natural way to interpret the 3D result is against neutrino-interaction
\emph{generators}. Because the published ancillary ships only binned 2D model
histograms, a quantitative 3D (or arbitrary-projection) comparison requires
generating truth-level events ourselves: running GENIE~\cite{Andreopoulos:2009rq}
(across tunes), NuWro~\cite{Golan:2012wx}, NEUT~\cite{Hayato:2009zz} and
GiBUU~\cite{Buss:2011mx} with the MINERvA ME~FHC flux, applying the same
fiducial phase space and the truth $\Eavail$ definition, and histogramming the
predicted $d^{3}\sigma$. No detector simulation is needed --- the unfolded
result already lives at truth level --- so this is a comparatively cheap,
high-value next step that directly exploits the unbinned measurement's ability
to discriminate models in the full $(\pT,\pPar,\Eavail)$ space.
